\sisection{S3. Fixed-pose transport across scoring functions}

Producers: \path{analysis/nonvina_scorer_transport.py}, with the ODDT stage carried by
\path{analysis/nonvina_scorer_oddt_worker.py}, and
\path{analysis/scorer_ordering_subset_stability.py} for S3.3. Frozen outputs are in
\path{results/nonvina_scorer_transport/} and
\path{results/scorer_ordering_subset_stability/}.

\sisubsection{S3.1 The pose block and the scoring implementations}

Every number in this section comes from one block of poses, 2{,}000 ligands by 9 DOCKSTRING
targets. Each of the 18{,}000 cells holds the top retained AutoDock Vina pose for that
ligand--target pair. The 9 targets are NR3C1, ADRB2, ADRB1, ADORA2A, AR, PTGS2, ACHE, DRD2
and LCK. Poses come from the released DOCKSTRING pose archives (Figshare DOI
10.6084/m9.figshare.16511577.v1), whose per-target MD5 is checked against a recorded value
at read time and whose extracted pose SDF carries its own SHA-256 in
\path{summary.json}. Ligands are a seed-20260802 subsample of the frozen seed-71
15{,}000-ligand DOCKSTRING support used elsewhere in this package. No ligand was dropped for
a missing pose. The pose is identical across all eight scoring outputs, so any difference
between two columns of the block comes from the scoring function.

smina 2020.12.10 (conda-forge build b08c07c, built 15 July 2025, based on AutoDock Vina
1.1.2) supplied two of the outputs, invoked in score-only mode with scoring arguments
\texttt{vina} and \texttt{vinardo} \cite{smina2013,vinardo2016}. The binary is identified in
\path{summary.json} by SHA-256 \texttt{b6902091\ldots}. smina does not echo SDF data tags in
score-only mode, so its scores are aligned to input record order, the per-target record count
is asserted, and the smina Vina column is cross-checked against the released DOCKSTRING Vina
column: per-target Pearson $r$ runs from 0.975 on LCK to 0.985 on NR3C1
(\path{summary.json}, \path{smina_vina_reproduction_check}).

ODDT 0.7 supplied five outputs: RF-Score v1, v2 and v3, NNScore 2.0 and PLECscore linear
\cite{oddt2015,rfscore2010,ballester2014rfscore,nnscore2011,plecscore2019}. RF-Score and
NNScore were fitted inside the run on the ODDT PDBbind 2016 descriptor tables, 3{,}767
training complexes, ODDT internal seed 1; on the 290-complex core set the fitted models reach
Pearson $r$ 0.788 for v1, 0.809 for v2, 0.797 for v3 and 0.819 for NNScore. PLECscore uses
the pretrained linear coefficients ODDT ships from that release
(\path{plecscore_linear_p5_l1_s65536_pdbbind2016.json}, SHA-256
\texttt{319a5041\ldots}); ODDT does not ship the matching PDBbind descriptor table, so no
offline hold-out check is possible for it. Receptors are the \texttt{dockstring} package
PDBQT resources, each recorded by SHA-256 and by protein atom count, from 1{,}852 atoms for
ADRB2 to 4{,}441 for PTGS2. The DRD2 receptor fails the default RDKit valence check, so a
relaxed sanitisation was applied uniformly to all nine receptors and verified bit-identical
to the default reader wherever the default reader succeeds.

The ODDT stage runs under a separate Python 3.9.23 interpreter with NumPy 1.26.4 and SciPy
1.13.1, because ODDT 0.7 does not install alongside the release Python 3.12 environment. Two
ODDT symbols, \path{oddt.scoring.sparse_vstack} and
\path{oddt.scoring.descriptors.sparse_vstack}, are patched at import because
\path{scipy.sparse.vstack} no longer accepts a map object; the iterator is materialised into
a list before stacking and the numeric effect is none. Reproducing this stage also needs about
600~MB of pose archives and an external smina binary, neither of which the release bundle
carries. The stage therefore sits outside \texttt{make all}, and its outputs are versioned and
checksum-verified instead of rebuilt: \path{fixed_pose_scores.csv.gz} at SHA-256
\texttt{8c0d10d2\ldots} and \path{scorer_spectral_metrics.csv} at SHA-256
\texttt{29ecc126\ldots}, both recorded in \path{output_checksums.json} and repeated inside
\path{summary.json}. The release-environment step reads those two files and recomputes every
statistic reported below.

Scores were not clipped at zero anywhere in this section. Clipping is a Vina-scale convention,
and applying it to the Vina-family columns alone would treat two of the eight outputs
differently from the rest; the released Vina column holds 3 positive values on this block. Each
column is multiplied by the polarity recorded in \path{scorer_spectral_metrics.csv}, $-1$ for
the Vina-form outputs and $+1$ for the learned outputs, so that larger values mean stronger
predicted binding. The PLECscore coefficients ODDT ships carry \texttt{intercept\_}~$=0$, so
absolute PLECscore values are offset from the $\mathrm{p}K_\mathrm{d}$ scale; every statistic
reported here is correlation-based and invariant to a per-column constant shift.

\sisubsection{S3.2 The four criteria and the per-scorer outcome}

Four criteria were applied to each of the eight outputs. The first asks for a raw PC1 variance
fraction of at least 0.50. The second asks for a cosine of at least 0.95 between the raw PC1
loading vector, signed so its entries sum to a positive number, and the uniform target
direction $\mathbf{1}/\sqrt{9}$, with every loading strictly positive. The third asks for a raw
participation ratio no larger than half the target count, 4.5 on nine targets. The fourth asks
that the lower end of the conservative 95\% cluster-bootstrap interval on the
centring-induced participation-ratio increase be strictly positive.

Intervals are the conservative union of two cluster bootstraps at 5{,}000 replicates each,
each resampling whole chemical groups with replacement. One scheme uses 1{,}957 Butina
clusters, built for this block by RDKit Butina clustering on Morgan radius-2 2048-bit
Tanimoto distances at cutoff 0.35, since DOCKSTRING carries no upstream cluster column
\cite{butina1999,morgan1965,rogers2010}; the other uses 1{,}876 Bemis--Murcko scaffold groups
\cite{murcko1996}. The two bootstrap seeds are 20260806 and 20260805. These are
finite-support sensitivity ranges on a fixed nine-target panel, not confidence intervals for a
superpopulation of targets.


The second criterion has two halves, and the producer records an interval for only one of
them. It bootstraps the cosine but reports the minimum PC1 loading as a point value, so the
requirement that every loading be positive carried no uncertainty. That gap matters for
RF-Score v2, whose minimum loading is 0.0469. We recomputed the missing interval from the
frozen per-cell scores with
\path{analysis/nonvina_scorer_min_loading_interval.py}, reusing the producer's own metric
function, clustering recipe and seeds, so no rescoring, pose archive or second interpreter is
needed. All eight intervals exclude zero. The narrowest margin is RF-Score v2 at
$[0.0275, 0.0653]$, and the widest loading is PLECscore at 0.2235 with interval
$[0.1978, 0.2474]$. The positivity half of the second criterion therefore holds for every
output under chemical resampling. The same run recomputes the cosine intervals as a check and
reproduces the released values to a maximum absolute difference of $1.8\times10^{-4}$; the
residual difference comes from recomputing the Butina and Bemis--Murcko labels rather than
reading the original assignment, so the resampled draws are not bit-identical.

\begin{table}[htbp]
\centering
\caption{Spectral summaries of the eight scoring outputs on the identical 2{,}000-by-9 fixed
pose block. Brackets give the conservative 95\% cluster-bootstrap interval, the union of the
Butina-cluster and Bemis--Murcko-scaffold schemes at 5{,}000 replicates each. The minimum
PC1 loading is reported without an interval in the producer's own table; its interval here comes
from the recomputation described in the text. Sources:
\protect\path{results/nonvina_scorer_transport/scorer_spectral_metrics.csv} and
\protect\path{results/nonvina_scorer_min_loading_interval/min_loading_intervals.csv}.}
\label{tab:s3metrics}
\small
\resizebox{\textwidth}{!}{%
\begin{tabular}{lllrlll}
\toprule
Output & PC1 fraction & Uniform cosine & Min.\ loading & Raw PR & Centred PR & PR increase \\
\midrule
Released Vina & 0.573 [0.552, 0.597] & 0.950 [0.939, 0.959] & 0.082 & 2.740 [2.562, 2.905] & 5.125 [4.930, 5.300] & 2.385 [2.157, 2.583] \\
smina Vina & 0.577 [0.557, 0.599] & 0.955 [0.945, 0.962] & 0.095 & 2.710 [2.548, 2.866] & 5.131 [4.937, 5.309] & 2.421 [2.197, 2.618] \\
Vinardo & 0.567 [0.551, 0.585] & 0.967 [0.953, 0.977] & 0.153 & 2.806 [2.673, 2.933] & 5.736 [5.497, 5.934] & 2.930 [2.658, 3.169] \\
RF-Score v1 & 0.635 [0.614, 0.657] & 0.959 [0.952, 0.964] & 0.064 & 2.345 [2.208, 2.484] & 6.444 [6.260, 6.580] & 4.099 [3.839, 4.306] \\
RF-Score v2 & 0.669 [0.648, 0.691] & 0.954 [0.947, 0.959] & 0.047 & 2.130 [2.008, 2.249] & 5.652 [5.470, 5.801] & 3.522 [3.283, 3.727] \\
RF-Score v3 & 0.630 [0.609, 0.652] & 0.956 [0.950, 0.962] & 0.057 & 2.369 [2.227, 2.510] & 5.999 [5.805, 6.155] & 3.630 [3.361, 3.857] \\
NNScore 2.0 & 0.526 [0.507, 0.545] & 0.973 [0.967, 0.977] & 0.134 & 3.210 [3.032, 3.397] & 6.795 [6.620, 6.922] & 3.585 [3.304, 3.814] \\
PLECscore linear & 0.414 [0.396, 0.432] & 0.986 [0.980, 0.989] & 0.223 & 4.593 [4.335, 4.848] & 7.465 [7.337, 7.532] & 2.872 [2.571, 3.118] \\
\bottomrule
\end{tabular}
}
\end{table}

\begin{table}[htbp]
\centering
\caption{Outcome of the four criteria for each output. Thresholds are 0.50 on the raw PC1
variance fraction, 0.95 on the uniform cosine with every PC1 loading strictly positive, 4.5 on
the raw participation ratio, and a strictly positive lower end of the conservative 95\%
cluster-bootstrap interval on the participation-ratio increase. Source:
\protect\path{results/nonvina_scorer_transport/summary.json}, \texttt{prediction\_outcome}.}
\label{tab:s3criteria}
\small
\begin{tabular}{lccccc}
\toprule
Output & PC1 $\geq 0.50$ & Cosine $\geq 0.95$ & Raw PR $\leq 4.5$ & PR rise $> 0$ & All four \\
\midrule
Released Vina & pass & pass & pass & pass & pass \\
smina Vina & pass & pass & pass & pass & pass \\
Vinardo & pass & pass & pass & pass & pass \\
RF-Score v1 & pass & pass & pass & pass & pass \\
RF-Score v2 & pass & pass & pass & pass & pass \\
RF-Score v3 & pass & pass & pass & pass & pass \\
NNScore 2.0 & pass & pass & pass & pass & pass \\
PLECscore linear & fail & pass & fail & pass & fail \\
\bottomrule
\end{tabular}
\end{table}

Seven outputs meet all four criteria. PLECscore fails two. Its PC1 fraction of 0.414 sits
below the 0.50 gate and its whole interval, [0.396, 0.432], stays below the gate. Its raw
participation ratio of 4.593 sits above the 4.5 gate, but the interval [4.335, 4.848] crosses
it, so that failure is the weaker of the two. PLECscore passes the other two criteria, and its
uniform cosine of 0.986 is the largest of the eight outputs: its leading component is the most
nearly uniform of any output while carrying the least variance. All eight pass the centring
criterion, with the smallest lower bound 2.157 on released Vina.

Two point estimates sit close to a gate on the other side. The released Vina cosine is 0.950
and the smina Vina cosine is 0.955, and both intervals reach below 0.95, to 0.939 and 0.945.
The criteria are evaluated on point estimates except for the fourth, which is defined on the
interval. The criteria, the code that evaluates them and the results they judge all enter the
repository in a single commit, so a reader of the released package cannot verify that the
criteria were fixed before the outputs were computed. Read them as a descriptive checklist
with stated thresholds.

\sisubsection{S3.3 Feature-class ordering under target subsetting}

The main text reports one ordering of the eight outputs by PC1 variance fraction: the smallest
of the three RF-Scores above the largest of the three empirical-term outputs, that above
NNScore, and NNScore above PLECscore. The frozen file states the rule as
\texttt{min(atom-pair counts) > max(empirical terms) > learned contact net > interaction
fingerprint}, evaluated on PC1 variance fraction.

We enumerated every target subset of every size from 3 to 9 on the same fixed pose block,
recomputed the eight PC1 fractions inside each subset, and recorded whether the ordering above
held there. The enumeration is exhaustive, with no sampling and no seed, so the counts in
Table~\ref{tab:s3subsets} are exact.

\begin{table}[htbp]
\centering
\caption{Exhaustive enumeration of the feature-class ordering under target subsetting on the
fixed nine-target pose block. Source:
\protect\path{results/scorer_ordering_subset_stability/subset_counts.csv}.}
\label{tab:s3subsets}
\small
\begin{tabular}{rrrr}
\toprule
Targets in subset & Subsets & Ordering holds & Fraction \\
\midrule
3 & 84 & 31 & 0.369 \\
4 & 126 & 76 & 0.603 \\
5 & 126 & 98 & 0.778 \\
6 & 84 & 78 & 0.929 \\
7 & 36 & 36 & 1.000 \\
8 & 9 & 9 & 1.000 \\
9 & 1 & 1 & 1.000 \\
\bottomrule
\end{tabular}
\end{table}

The ordering holds on the full panel, on all 9 leave-one-target-out panels and on all 36
seven-target panels. It breaks on 6 of the 84 six-target panels and on 53 of the 84
three-target panels. What the enumeration bounds is how far the full-panel ordering can be
read as general: it is a property of this nine-target panel that survives removal of two
targets and degrades from there. It carries nothing to targets outside the panel or to scoring
functions outside the eight. It also does not simulate the poses another function would have
selected in its own search.

\sisubsection{S3.4 Scorer lineage dependence}

The eight outputs are not eight independent scoring functions. They fall into four lineages.

The first is the AutoDock Vina empirical sum, present twice: the released DOCKSTRING column
and smina's reimplementation of the same function, which agree at per-target Pearson $r$
between 0.975 and 0.985. The smina Vina column is present only as a same-program control, so
that a Vinardo-versus-Vina difference cannot be charged to the change of executable. Vinardo
belongs to
the same empirical form and is evaluated by the same smina binary. It changes three things in
that form: the gauss2 term is dropped, repulsion is reweighted, and the hydrophobic and
hydrogen-bond terms are redefined.

The second lineage is RF-Score, present in three versions of one model. Version 1 is a random
forest over 36 element-pair close-contact counts within 12~\AA{}. Version 2 bins the same
contacts over multiple distance shells. Version 3 takes the version-1 contacts plus the
AutoDock Vina intermolecular terms as extra input features, so it is not independent of the
Vina function and is reported as a graded intermediate rather than a separate lineage.

The third lineage is NNScore 2.0, an ensemble of neural networks over 350 BINANA descriptors.
The fourth is PLECscore, a linear model over the protein--ligand extended connectivity
fingerprint. Eight outputs map onto four lineages, and two of the four are linked through
RF-Score v3, whose input includes the Vina intermolecular terms.

Agreement of the per-ligand shared axis with the released Vina column tracks lineage. The
Spearman correlation of the leading ligand-side component with the released Vina ligand axis is
0.962 for smina Vina, 0.757 for Vinardo, 0.507 for RF-Score v2, 0.506 for RF-Score v3, 0.425
for RF-Score v1, 0.368 for NNScore and 0.317 for PLECscore
(\path{summary.json}, \texttt{shared\_axis\_agreement}). The same axis correlates with ligand
heavy-atom count at Spearman 0.367 for PLECscore up to 0.706 for RF-Score v3.

Two limits bound what this block can show. No commercial scoring function was tested; Glide and
GOLD are absent. And rescoring retained poses is not re-docking: each function is evaluated on
the pose AutoDock Vina selected, so a function that would have selected a different pose in its
own search is not simulated. Agreement could therefore differ in either direction in a full
per-scorer re-docking run. RF-Score, NNScore and PLECscore were trained on
PDBbind crystal complexes and are applied off-distribution to docked poses. Nothing in this
block shows that any scorer is accurate. Nothing in it shows that the shared axis is
biologically meaningful, or that the residual carries target-resolved mechanism.
